\section{CSV Data Importing}

\begin{frame}[fragile]{What Is CSV Data Importing?}
	\begin{itemize}
		\item CSV files can also load records during module install/update.
		\item Best for tabular, repetitive data such as:
		\begin{itemize}
			\item access rights (\texttt{ir.model.access.csv})
			\item simple master lists
			\item bulk demo rows
		\end{itemize}
\begin{block}{Manifest Entry}
\begin{lstlisting}
'data': [
	'security/ir.model.access.csv',
	'data/student_level_data.csv',
],
\end{lstlisting}
\end{block}
	\end{itemize}
\end{frame}

\begin{frame}[fragile]{CSV File Structure}
\begin{lstlisting}
id,name,code,active
student_level_beginner,Beginner,BEG,True
student_level_intermediate,Intermediate,INT,True
student_level_advanced,Advanced,ADV,True
\end{lstlisting}
	\begin{itemize}
		\item First row = field names.
		\item \texttt{id} column = XML / external ID (without module prefix in file).
		\item Other columns map to model fields.
	\end{itemize}
\end{frame}

\begin{frame}[fragile]{How Odoo Knows the Model}
	\begin{itemize}
		\item For most CSV files, the model is deduced from the filename / import conventions used by the module loader.
		\item Classic and most common case:
\begin{lstlisting}
security/ir.model.access.csv
\end{lstlisting}
		\item Access rights CSV targets \texttt{ir.model.access}.
	\end{itemize}
\end{frame}

\begin{frame}[fragile]{Access Rights CSV Example}
\begin{lstlisting}
id,name,model_id:id,group_id:id,perm_read,perm_write,perm_create,perm_unlink
access_student_user,student.user,model_student_student,base.group_user,1,0,0,0
access_student_manager,student.manager,model_student_student,school.group_school_manager,1,1,1,1
\end{lstlisting}
	\begin{itemize}
		\item \texttt{model\_id:id} and \texttt{group\_id:id} use XML IDs.
		\item \texttt{1}/\texttt{0} mean True/False for permissions.
	\end{itemize}
\end{frame}

\begin{frame}[fragile]{Relational Values in CSV}
\begin{lstlisting}
id,name,classroom_id/id,gender
student_ahmed,Ahmed Muhumed,classroom_grade_12,male
student_amina,Amina Ali,classroom_grade_12,female
\end{lstlisting}
	\begin{itemize}
		\item \texttt{classroom\_id/id} writes a many2one through an XML ID.
		\item This is cleaner than numeric database IDs.
	\end{itemize}
\end{frame}

\begin{frame}[fragile]{XML vs CSV}
	\begin{itemize}
		\item Use \textbf{XML} when you need:
		\begin{itemize}
			\item complex structures
			\item \texttt{noupdate}, \texttt{function}, rich field options
			\item menus/views/templates
		\end{itemize}
		\item Use \textbf{CSV} when you need:
		\begin{itemize}
			\item simple bulk rows
			\item access rights tables
			\item easy spreadsheet editing
		\end{itemize}
	\end{itemize}
\end{frame}

\begin{frame}{CSV Best Practices}
	\begin{itemize}
		\item Keep headers exact field names.
		\item Always provide an \texttt{id} column for stable external IDs.
		\item Prefer \texttt{field/id} for many2one relations.
		\item Validate CSV in a clean database install.
		\item Do not mix too much complex logic into CSV; switch to XML when needed.
	\end{itemize}
\end{frame}
